\section{Fathers' Jobs}\label{sec:fathers}

Parents' employers and networks shape where young people find their first job and raise their initial earnings
\citep{corak2011intergenerational, kramarz2014strong, staiger2022}. The PSID does not identify employers, and neither
do the public NLSY files for fathers, so working for the father's employer cannot be observed; the father's occupation
group and his occupation-by-industry cell, his modal codes when the son was 12--18, are the proxies.\footnote{The
father is the male head of the son's family at 16, so these measures exist for 52 percent of Black and 87 percent of
White sons. Fathers are coded in 1970 Census codes and sons after 2001 in 2000 or 2010 codes; the 35--55 window falls
more often after 2001, so coding drift lowers measured following after 35, which works against the convergence below.
Code: \texttt{32.IndustryCareers.py}, \texttt{35.EarlyIndustryLatePayoff.py} and
\texttt{43.FatherJobSelfEmployment.py}.}

\paragraph{Following pays at entry, not in growth.} Sons work in their father's industry group two to three times as
often as random matching within race and parental-income tercile would predict (14.2 percent of White sons at 35--55
against a benchmark of 5.1 percent, and 10.1 percent of Black sons against 5.4 percent). Holding the son's education
and the level and slope of his industry fixed, sons in their father's industry group earn 0.079$^{***}$ more at 25--34
and 0.054$^{*}$ more at 35--55, but their wages grow no faster (0.005, s.e.\ 0.005), and following accounts for none of
the racial divergence ($-0.002$ of 0.107$^{**}$).\footnote{Appendix Table~\ref{tab:fatherindustry} and
\texttt{output/industry\_careers.md}, Section 4. Few Black fathers worked in high-paying service industries (21 in
finance, insurance or real estate), so industry-specific following cannot be estimated for Black sons.} Family networks,
as measured here, transmit a starting advantage that is the same at every age at which the IGE is measured.

\paragraph{Convergence on the father's occupation.} Over the career, sons of richer parents converge on their fathers'
occupations (Table~\ref{tab:fatherse}). The parental-income gradient of being in the father's occupation group is
0.051$^{***}$ at 25--34 and 0.099$^{***}$ at 35--55 (change 0.052$^{***}$). The change is precisely estimated for White
sons (0.060$^{***}$) and imprecise for Black sons (0.036, s.e.\ 0.028); the two cannot be distinguished. The level of
following differs by race at every age: at 35--55, 22.4 percent of White and 15.3 percent of Black sons work in their
father's occupation group. Within the top parental tercile, sons of White parents are 11.7 percentage points more likely
than sons of Black parents to work in their father's occupation group (0.117$^{***}$, resting on 121 Black sons in that
tercile), while rich and poor White sons do not differ ($-0.086$, s.e.\ 0.061): the contrast is a racial gap in
following the father's occupation at equal parental income, not class amplifying it.

\paragraph{The father's exact job is a family business.} Whether a son holds his father's exact job---his occupation
group in his industry group---does not depend on parental income at 25--34 (0.004) and does at 35--55 (0.031$^{**}$;
change 0.023$^{**}$), because sons of richer parents who hold it in their early thirties keep it: among the 137 sons in
the father's cell at 25--34, 84 are still in it at 35--55 and 53 leave, and leaving falls with parental income
($-0.206^{***}$); entry after 34 is rare (39 of 1,892 sons) and unrelated to parental income. Only 9 of the 84 who keep
the father's job are Black. The self-employment split shows what the exact job is. Twenty-one percent of White and 8
percent of Black sons have a self-employed father, and the convergence comes entirely from them: among sons of
self-employed fathers the gradient rises by 0.065$^{***}$ after 35, among other sons by 0.009 (0.011), and among sons
for whom neither father nor son is self-employed by 0.005. The cells most often inherited after 35 are farming in
agriculture (31 sons, 56 percent of them self-employed and 75 percent with a self-employed father), craft work in
construction and management in construction. Convergence on the father's exact job is the inheritance of a farm or a
firm \citep[the business and wealth transmission of][]{fagereng2021wealthy}, not of a position in a network. The
occupation-level result is not: excluding self-employed fathers, the gradient still rises after 35 (0.037$^{*}$) and the
rich-White-versus-rich-Black contrast is unchanged (0.117$^{***}$). Keeping the father's job also does not depend on
self-employment ($-0.183^{**}$ among the 90 sons of fathers who were not self-employed), and holding the father's exact
job brings no wage premium at 35--55 and no faster wage growth.

\begin{table}[!htbp]
\centering
\caption{Fathers' jobs and self-employment (PSID, sons)}
\label{tab:fatherse}
\resizebox{\ifdim\width>\textwidth\textwidth\else\width\fi}{!}{%
\resizebox*{!}{\ifdim\totalheight>0.85\textheight 0.85\textheight\else\totalheight\fi}{%
\begin{tabular}{lcccc}
\toprule
 & \begin{tabular}[b]{@{}c@{}}All sons\end{tabular} & \begin{tabular}[b]{@{}c@{}}Father not\\self-employed\end{tabular} & \begin{tabular}[b]{@{}c@{}}Father\\self-employed\end{tabular} & \begin{tabular}[b]{@{}c@{}}Neither father\\nor son\\self-employed\end{tabular} \\
 & (1) & (2) & (3) & (4) \\
\midrule
\multicolumn{5}{l}{\textit{Son in the father's occupation $\times$ industry cell}} \\
All: log parental income, 25--34 & $0.004$ & $-0.006$ & $0.009$ & $-0.005$ \\
 & (0.012) & (0.012) & (0.023) & (0.013) \\
All: log parental income, 35--55 & $0.031^{**}$ & $0.009$ & $0.072^{***}$ & $0.008$ \\
 & (0.016) & (0.016) & (0.025) & (0.018) \\
All: log parental income $\times$ (age $\geq$ 35) & $0.023^{**}$ & $0.009$ & $0.065^{***}$ & $0.005$ \\
 & (0.010) & (0.011) & (0.021) & (0.012) \\
White: log parental income, 25--34 & $-0.005$ & $-0.011$ & $-0.005$ & $-0.007$ \\
 & (0.014) & (0.014) & (0.025) & (0.015) \\
White: log parental income, 35--55 & $0.025$ & $0.006$ & $0.065^{**}$ & $0.006$ \\
 & (0.020) & (0.020) & (0.025) & (0.023) \\
White: log parental income $\times$ (age $\geq$ 35) & $0.026^{**}$ & $0.009$ & $0.073^{***}$ & $0.004$ \\
 & (0.012) & (0.012) & (0.022) & (0.015) \\
Son-years / sons at 35--55 & 13,835 / 2,166 & 10,539 / 1,683 & 2,399 / 373 & 9,324 / 1,579 \\
\addlinespace
\multicolumn{5}{l}{\textit{Son in the father's occupation group}} \\
All: log parental income, 25--34 & $0.051^{***}$ & $0.040^{**}$ & $0.067^{**}$ & $0.037^{*}$ \\
 & (0.016) & (0.019) & (0.030) & (0.019) \\
All: log parental income, 35--55 & $0.099^{***}$ & $0.071^{***}$ & $0.163^{***}$ & $0.072^{***}$ \\
 & (0.020) & (0.026) & (0.027) & (0.027) \\
All: log parental income $\times$ (age $\geq$ 35) & $0.052^{***}$ & $0.037^{*}$ & $0.104^{***}$ & $0.038^{*}$ \\
 & (0.016) & (0.021) & (0.026) & (0.022) \\
White: log parental income, 25--34 & $0.040^{**}$ & $0.035$ & $0.042$ & $0.033$ \\
 & (0.018) & (0.022) & (0.032) & (0.022) \\
White: log parental income, 35--55 & $0.099^{***}$ & $0.073^{**}$ & $0.161^{***}$ & $0.076^{**}$ \\
 & (0.023) & (0.029) & (0.027) & (0.029) \\
White: log parental income $\times$ (age $\geq$ 35) & $0.060^{***}$ & $0.039^{*}$ & $0.124^{***}$ & $0.041^{*}$ \\
 & (0.017) & (0.022) & (0.028) & (0.024) \\
Son-years / sons at 35--55 & 13,908 / 2,174 & 10,604 / 1,691 & 2,404 / 373 & 9,383 / 1,588 \\
\addlinespace
\multicolumn{5}{l}{\textit{Race, class and keeping the father's job}} \\
Rich White $-$ rich Black, father's occupation group, 35--55 & $0.117^{***}$ & $0.117^{***}$ & $-0.019$ & $0.113^{***}$ \\
 & (0.038) & (0.040) & (0.143) & (0.043) \\
Rich White $-$ rich Black, father's cell, 35--55 & $0.011$ & $-0.008$ & $0.131^{***}$ & $-0.004$ \\
 & (0.029) & (0.030) & (0.033) & (0.033) \\
Black sons in the top tercile & 121 & 112 & 9 & 105 \\
Leaving the father's cell (in at 25--34) on log parental income & $-0.206^{***}$ & $-0.183^{**}$ & $-0.223^{*}$ & $-0.163^{**}$ \\
 & (0.066) & (0.080) & (0.126) & (0.079) \\
Sons in the father's cell at 25--34 & 137 & 90 & 46 & 82 \\
\bottomrule
\end{tabular}}}
\par\smallskip
\begin{minipage}{\textwidth}\footnotesize \textit{Notes:} Son-year linear probability models, Lee--Solon controls (quartics in the parent's age and the son's age), clusters 1968 family; all sons weighted, White sons unweighted. The father's cell is his modal 11 occupation groups $\times$ 33 industry groups when the son is 12--18; the father is self-employed when at least half of his valid interviews from 1976 on, with the son aged 10--25, report working for himself (alone or also for someone else). The change row is from one regression on 25--55 with log parental income $\times$ (age $\geq$ 35). Rich White $-$ rich Black: sons in the top pooled tercile of parental income, race $\times$ tercile cells with log parental income, weighted. Leaving: sons in the father's cell in their modal 25--34 job, out of it in their modal 35--55 job, on log parental income and a quadratic in birth cohort. These estimates use sons with a father's self-employment code (and, for leaving, a cohort quadratic and no weights); on all sons with a father's cell the same contrasts are 0.131$^{***}$ (0.038) for rich White $-$ rich Black and $-0.167^{**}$ (0.067) for leaving (	exttt{35.EarlyIndustryLatePayoff.py}, Section F). $^{***}$ $p<0.01$, $^{**}$ $p<0.05$, $^{*}$ $p<0.1$. Source: \texttt{43.FatherJobSelfEmployment.py} (\texttt{output/father\_job\_selfemp.csv}).\end{minipage}
\end{table}


The fathers' evidence adds two channels to the career evidence: a level advantage at entry through family networks,
which accounts for $-0.002$ (0.004) of the 0.107$^{**}$ change in the racial gap in wage persistence on the father-present
sample, and the inheritance of family businesses, a matter of assets rather than access, which accounts for $-0.004$ of a
0.120$^{**}$ change. Convergence on fathers' occupations outside the family firm, with its racial gap at equal parental
income, is a third. None of them carries the growth of wages that produces the divergence; that is the job at 25--34,
its access, and schooling.
